\section{Solmut ja algebralliset käyrät}

Matematiikassa solmuja tarkastellaan ympyrännäköisinä topologisina olioina, jotka saavat sotkeutua itseensä. Epämuodollisesti määriteltynä matemaattinen solmu on kuin tavallinen solmu, jonka päät on liimattu yhteen. Linkillä tarkoitetaan kappaletta, joka koostuu erillisistä solmuista, jotka voivat olla kietoutuneina toisiinsa. Työssä noudatetaan Roflsenin kirjan \textit{Knots and Links} merkintätapoja.
\begin{definition}
    Topologisen avaruuden $X$ osajoukko $K\subset X$ on \textit{solmu}, jos $K$ on homeomorfinen palloon $\pallo^p$, jollekin $p\in\N$. $K$ on \textit{linkki}, jos se on homeomorfinen useiden pallojen erilliseen yhdisteeseen $\bigsqcup_{i=1}^n \pallo^p_i$.
\end{definition}

